% =============================================================
% Bibliografía
% =============================================================
% Formato APA según indica la plantilla de la EPSG.
% De momento se usa thebibliography manual; si conviene se migra a
% bibtex/biblatex con un archivo .bib más adelante.
% =============================================================

\begin{thebibliography}{9}

% --- Algoritmos genéticos / computación evolutiva ---

\bibitem{holland1975}
Holland, J. H. (1975). \textit{Adaptation in Natural and Artificial Systems}. University of Michigan Press.

\bibitem{goldberg1989}
Goldberg, D. E. (1989). \textit{Genetic Algorithms in Search, Optimization, and Machine Learning}. Addison-Wesley.

\bibitem{eiben2015}
Eiben, A. E., \& Smith, J. E. (2015). \textit{Introduction to Evolutionary Computing} (2.\textsuperscript{a} ed.). Springer.

\bibitem{mitchell1996}
Mitchell, M. (1996). \textit{An Introduction to Genetic Algorithms}. MIT Press.

\bibitem{jinbranke2005}
Jin, Y., \& Branke, J. (2005). Evolutionary optimization in uncertain environments: A survey. \textit{IEEE Transactions on Evolutionary Computation, 9}(3), 303-317.

\bibitem{blum2003}
Blum, C., \& Roli, A. (2003). Metaheuristics in combinatorial optimization: Overview and conceptual comparison. \textit{ACM Computing Surveys, 35}(3), 268-308.

\bibitem{millergoldberg1995}
Miller, B. L., \& Goldberg, D. E. (1995). Genetic algorithms, tournament selection, and the effects of noise. \textit{Complex Systems, 9}(3), 193-212.

\bibitem{eiben1999}
Eiben, A. E., Hinterding, R., \& Michalewicz, Z. (1999). Parameter control in evolutionary algorithms. \textit{IEEE Transactions on Evolutionary Computation, 3}(2), 124-141.

\bibitem{crepinsek2013}
Črepinšek, M., Liu, S.-H., \& Mernik, M. (2013). Exploration and exploitation in evolutionary algorithms: A survey. \textit{ACM Computing Surveys, 45}(3), artículo 35.

\bibitem{mahfoud1995}
Mahfoud, S. W. (1995). \textit{Niching methods for genetic algorithms} [Tesis doctoral, University of Illinois at Urbana-Champaign].

\bibitem{jin2011}
Jin, Y. (2011). Surrogate-assisted evolutionary computation: Recent advances and future challenges. \textit{Swarm and Evolutionary Computation, 1}(2), 61-70.

\bibitem{mouret2015}
Mouret, J.-B., \& Clune, J. (2015). \textit{Illuminating search spaces by mapping elites}. arXiv:1504.04909.

\bibitem{pugh2016}
Pugh, J. K., Soros, L. B., \& Stanley, K. O. (2016). Quality diversity: A new frontier for evolutionary computation. \textit{Frontiers in Robotics and AI, 3}, artículo 40.

% --- Computación evolutiva e inteligencia artificial en juegos ---

\bibitem{yannakakis2018}
Yannakakis, G. N., \& Togelius, J. (2018). \textit{Artificial Intelligence and Games}. Springer.

\bibitem{risi2017}
Risi, S., \& Togelius, J. (2017). Neuroevolution in games: State of the art and open challenges. \textit{IEEE Transactions on Computational Intelligence and AI in Games, 9}(1), 25-41.

\bibitem{togelius2011}
Togelius, J., Yannakakis, G. N., Stanley, K. O., \& Browne, C. (2011). Search-based procedural content generation: A taxonomy and survey. \textit{IEEE Transactions on Computational Intelligence and AI in Games, 3}(3), 172-186.

% --- Magic: The Gathering ---

\bibitem{guinness}
Guinness World Records. (s.f.). \textit{Most played trading card game}. Recuperado de \url{https://www.guinnessworldrecords.com/world-records/most-played-trading-card-game}

\bibitem{wotc_standard}
Wizards of the Coast. (s.f.). \textit{Standard: reglas y rotación del formato}. Recuperado de \url{https://magic.wizards.com/en/formats/standard}

\bibitem{wotc_rules}
Wizards of the Coast. (2025). \textit{Magic: The Gathering Comprehensive Rules}. Recuperado de \url{https://magic.wizards.com/en/rules}

\bibitem{karsten_manabase}
Karsten, F. (2018). \textit{How many sources do you need to consistently cast your spells?} ChannelFireball. Recuperado de \url{https://www.channelfireball.com/}

\bibitem{churchill2019}
Churchill, A., Biderman, S., \& Herrick, A. (2019). \textit{Magic: The Gathering is Turing Complete}. arXiv:1904.09828. Recuperado de \url{https://arxiv.org/abs/1904.09828}

% --- Herramientas ---

\bibitem{forge}
Card-Forge contributors. \textit{Forge: an unofficial Magic: The Gathering rules engine} [Software]. GitHub. Recuperado de \url{https://github.com/Card-Forge/forge}

\bibitem{scryfall}
Scryfall. \textit{Scryfall API Documentation}. Recuperado de \url{https://scryfall.com/docs/api}

% --- Estado del arte ---

\bibitem{bjorke2017}
Bjørke, S. J., \& Fludal, K. A. (2017). \textit{Deckbuilding in Magic: The Gathering using a genetic algorithm} [Tesis de máster, Norwegian University of Science and Technology]. Department of Computer Science, NTNU.

\bibitem{mira2021}
Mira Abad, A. (2021). \textit{Aprendizaje de estrategias de juego en juegos de cartas complejos y estratégicos mediante técnicas de inteligencia artificial} [Trabajo Final de Grado]. Escola Politècnica Superior de Gandia, Universitat Politècnica de València.

\bibitem{cowling2012}
Cowling, P. I., Ward, C. D., \& Powley, E. J. (2012). Ensemble determinization in Monte Carlo tree search for the imperfect information card game Magic: The Gathering. \textit{IEEE Transactions on Computational Intelligence and AI in Games, 4}(4), 241-257.

\bibitem{zilio2018}
Zilio, F., Prates, M. O. R., \& Lamb, L. C. (2018). \textit{Neural networks models for analyzing Magic: The Gathering cards}. arXiv:1810.03744. Presentado en ICONIP 2018.

\bibitem{garciasanchez2016}
García-Sánchez, P., Tonda, A., Squillero, G., Mora, A. M., \& Merelo, J. J. (2016). Evolutionary deckbuilding in HearthStone. En \textit{2016 IEEE Conference on Computational Intelligence and Games (CIG)} (pp. 1-8). IEEE.

\bibitem{mahlmann2012}
Mahlmann, T., Togelius, J., \& Yannakakis, G. N. (2012). Evolving card sets towards balancing Dominion. En \textit{2012 IEEE Congress on Evolutionary Computation (CEC)} (pp. 1-8). IEEE. \url{https://doi.org/10.1109/CEC.2012.6256441}

\bibitem{miles2006}
Miles, C., \& Louis, S. J. (2006). Towards the co-evolution of influence map tree based strategy game players. En \textit{2006 IEEE Symposium on Computational Intelligence and Games (CIG)} (pp. 75-82). IEEE. \url{https://doi.org/10.1109/CIG.2006.311684}

\bibitem{garciasanchez2018}
García-Sánchez, P., Tonda, A., Mora, A. M., Squillero, G., \& Merelo, J. J. (2018). Automated playtesting in collectible card games using evolutionary algorithms: A case study in HearthStone. \textit{Knowledge-Based Systems, 153}, 133-146.

\bibitem{kowalski2020}
Kowalski, J., \& Miernik, R. (2020). Evolutionary approach to collectible card game arena deckbuilding using active genes. En \textit{2020 IEEE Congress on Evolutionary Computation (CEC)} (pp. 1-8). IEEE. \url{https://arxiv.org/abs/2001.01326}

\bibitem{miernik2022}
Miernik, R., \& Kowalski, J. (2022). Evolving evaluation functions for collectible card game AI. En \textit{Proceedings of the 14th International Conference on Agents and Artificial Intelligence (ICAART)} (Vol. 3, pp. 253-260). SciTePress. \url{https://arxiv.org/abs/2105.01115}

\bibitem{mesentier2019}
de Mesentier Silva, F., Canaan, R., Lee, S., Fontaine, M. C., Togelius, J., \& Hoover, A. K. (2019). Evolving the hearthstone meta. En \textit{2019 IEEE Conference on Games (CoG)} (pp. 1-8). IEEE. \url{https://doi.org/10.1109/CIG.2019.8847966}

\bibitem{fontaine2019}
Fontaine, M. C., Lee, S., Soros, L. B., de Mesentier Silva, F., Togelius, J., \& Hoover, A. K. (2019). Mapping hearthstone deck spaces through MAP-Elites with sliding boundaries. En \textit{Proceedings of the Genetic and Evolutionary Computation Conference (GECCO)} (pp. 161-169). ACM.

\bibitem{zhang2022}
Zhang, Y., Fontaine, M. C., Hoover, A. K., \& Nikolaidis, S. (2022). Deep surrogate assisted MAP-Elites for automated HearthStone deckbuilding. En \textit{Proceedings of the Genetic and Evolutionary Computation Conference (GECCO)} (pp. 158-167). ACM.

\bibitem{stiegler2016}
Stiegler, A., Messerschmidt, C., Maucher, J., \& Dahal, K. (2016). Hearthstone deck-construction with a utility system. En \textit{2016 10th International Conference on Software, Knowledge, Information Management \& Applications (SKIMA)}. IEEE.

\bibitem{bhatt2018}
Bhatt, A., Lee, S., de Mesentier Silva, F., Watson, C. W., Togelius, J., \& Hoover, A. K. (2018). Exploring the hearthstone deck space. En \textit{Proceedings of the 13th International Conference on the Foundations of Digital Games (FDG)}. ACM.

\bibitem{vieira2020}
Vieira, R. E. S., Tavares, A. R., \& Chaimowicz, L. (2020). Drafting in collectible card games via reinforcement learning. En \textit{Proceedings of the 19th Brazilian Symposium on Computer Games and Digital Entertainment (SBGames)}. IEEE.

\bibitem{swiechowski2018}
Świechowski, M., Tajmajer, T., \& Janusz, A. (2018). Improving Hearthstone AI by combining MCTS and supervised learning algorithms. En \textit{2018 IEEE Conference on Computational Intelligence and Games (CIG)}. IEEE.

\end{thebibliography}
